\chapter{Packages: Reuse, Privacy, and the Build}
\label{ch:packages}\index{package}

The previous chapter introduced the package in passing -- enough to read the code.
This chapter argues that the package is the single most important structural idea in
Ada, and the one that shapes this whole codebase: one peripheral per package, a tidy
\texttt{ESP32S3.*} hierarchy, a private type guarding every hardware resource, and a
build that recompiles exactly what changed. The payoff is worth stating up front,
because it is unusual. In Ada \emph{three} things that are separate, manually-managed
mechanisms in C collapse into one construct:

\begin{itemize}
\item the unit of \textbf{privacy} -- what is hidden from the rest of the program;
\item the unit of \textbf{separate compilation} -- what recompiles together;
\item the unit of \textbf{dependency} -- what depends on what.
\end{itemize}

\noindent
In C those are headers, translation units, and link order: three things that can
disagree. Folding them into the package, with a compiler-enforced split between
interface and implementation, is what gives Ada strong encapsulation, fast and
\emph{correct} incremental builds, and collision-free naming at scale, all at once.

\section{A contract, split from its implementation}\index{package!specification}\index{package!body}

A package comes in two files: a \emph{specification} (\texttt{.ads}) and a
\emph{body} (\texttt{.adb}). The spec is the contract -- everything a client may
rely on. The body is the implementation -- which nobody else can see.

\begin{lstlisting}
--  esp32s3-uart.ads  -- the SPEC: all a client is allowed to know
package ESP32S3.UART is
   type Session is limited private;             --  a handle; its insides are hidden
   procedure Acquire (S : in out Session; Port : UART_Port);
   procedure Write   (S : Session; Data : Byte_Array);
private
   type Session is limited record ... end record; --  visible to the compiler, not to clients
end ESP32S3.UART;
\end{lstlisting}

\begin{lstlisting}
--  esp32s3-uart.adb  -- the BODY: the implementation; sealed
package body ESP32S3.UART is
   procedure Write (S : Session; Data : Byte_Array) is ... end Write;
end ESP32S3.UART;
\end{lstlisting}

A caller is compiled against the spec \emph{alone}. It cannot name, reach, or depend
on anything in the body. This single rule is the root of almost everything that
follows: because clients see only the spec, the body is free to change underneath
them without breaking -- or even recompiling -- a single line of client code.

\cnote{A C header is just text the preprocessor pastes in; nothing makes it agree
with the \texttt{.c} that implements it, so the two can silently drift, and a
\texttt{struct} declared in a header leaks its layout to everyone who includes it.
An Ada spec \emph{is} the interface, checked by the compiler against its body, and a
\texttt{private} type hands out a handle whose representation no client can see.}

\section{Reuse and separate compilation}\index{separate compilation}

The package is the unit of reuse. \texttt{ESP32S3.SPI} is written once and
\texttt{with}-ed by every client that needs it -- the W5500 Ethernet driver, the
W25Q flash, the LCD -- with no copy-paste and no textual inclusion. Each package is
compiled independently to an object plus an \texttt{.ali} (Ada Library Information)
file that records its interface; clients link against the result rather than
recompiling it.

Because the spec is the reuse boundary, you can replace a body wholesale and every
client still works untouched. This project does exactly that: the heap body was
swapped for a TLSF allocator (Chapter~\ref{ch:heap}) and every caller of
\texttt{new} and \texttt{malloc} was none the wiser, because the \emph{spec} did not
move.

\emph{Generics} extend reuse to parameterised packages -- a package (or subprogram)
templated over types and values, instantiated per use. The block-device and
container code use them to write an algorithm once and specialise it many times
without giving up type safety.

\section{Namespaces: no collisions by construction}\index{naming collision}\index{package!hierarchy}

Every declaration lives inside a package, so its full name is \texttt{Package.Name}.
Two packages may both export \texttt{Read} or \texttt{Reset} with no conflict
whatever: \texttt{ESP32S3.SPI.Read} and \texttt{ESP32S3.UART.Read} are simply
different names. There is no global namespace to pollute.

\cnote{This is the problem C solves with manual prefixes -- \texttt{uart\_read()},
\texttt{spi\_read()} -- because every external symbol shares one flat namespace.
Ada gives you the prefix for free, and it means something: it is the package.}

The hierarchy carries structure. \texttt{ESP32S3.W5500.Sockets} is a \emph{child} of
\texttt{ESP32S3.W5500}, itself a child of \texttt{ESP32S3}; the dotted name tells you
at a glance exactly where a name comes from, and the tree groups a subsystem without
one enormous file. Bringing names into scope is opt-in and granular:

\begin{lstlisting}
use ESP32S3.GPIO;                 --  drop the prefix for everything in GPIO
use type ESP32S3.UART.Byte;       --  bring in JUST one type's operators
package U renames ESP32S3.UART;   --  a short local alias for a long name
\end{lstlisting}

\noindent
And collisions are \emph{caught}, not guessed: if two \texttt{use}d packages both
export \texttt{Read}, an unqualified \texttt{Read} becomes illegal and the compiler
makes you say which one you mean -- rather than silently picking, as a C macro or a
flat symbol table would. Overloading (several subprograms of the same name, told
apart by parameter and result profile) composes with all of this: it resolves
\emph{within} a package, orthogonally to the qualification \emph{between} packages.

\section{Information hiding}\index{encapsulation}

This is where packages earn their keep on a system that has to be trusted. A private
type hands clients a handle they can pass around but whose representation they cannot
see or alter -- and the compiler enforces it; there is no \texttt{\&session->field}
escape hatch. Better still, anything declared in the \emph{body} cannot even be
\emph{named} from outside. This codebase's UART driver hides the raw register bus in
a package nested inside the body, reachable only through one ownership-checked
gateway:

\begin{lstlisting}
package body ESP32S3.UART is
   package State is
      function Owned (S : Session) return Bus;  --  raises unless S holds the port
   end State;
   --  the raw per-port Bus array lives in State's body -- nothing else can name it
end ESP32S3.UART;
\end{lstlisting}

\noindent
A transfer physically cannot reach the registers without proving ownership, because
there is no other door. The guard is impossible to forget or bypass -- not by
convention, but because the alternative does not exist in any scope a client can
reach. Layer on \texttt{limited private} plus a \texttt{controlled} type whose
\texttt{Finalize} releases the bus on scope exit, and the invariant ``the bus is
always released'' belongs to the package, not to every caller remembering to clean
up. The mechanics of private types, \texttt{controlled} finalisation and dispatching
are the subject of Chapter~\ref{ch:encap}; the point here is that the \emph{package}
is the boundary they all hang on.

\section{Child packages: extend without modifying}\index{package!child}

A \emph{public} child adds capability to a subsystem without touching or recompiling
its parent. The W5500 driver ships a core (\texttt{ESP32S3.W5500}) and layers the
socket engine, DNS and DHCP on top as children -- \texttt{ESP32S3.W5500.Sockets} and
friends -- each a separate file, none of which forces the parent to change. A child's
spec and body may see the parent's \emph{private} part, so a family of packages can
share one hidden representation while still hiding it from the world. A \emph{private}
child (\texttt{ESP32S3.Foo.Internal}) goes further: it is visible only \emph{within}
its subsystem, an implementation helper the outside literally cannot \texttt{with} --
finer-grained than C's file-static, because it scopes across several files of one
subsystem at once.

\section{How the build system uses packages}\index{gprbuild}\index{build system}

The package model is not just for humans; it drives compilation directly.

\textbf{Dependency is semantic, not textual.} \texttt{with ESP32S3.SPI;} records that
this unit depends on the \emph{interface} of \texttt{ESP32S3.SPI}. GNAT tracks these
relationships in the \texttt{.ali} files and builds a dependency graph from them.

\textbf{Recompilation is minimal and exact.} Change a package \emph{body} and only
that body recompiles -- its clients do not, because their view (the spec) has not
moved. Change a \emph{spec} and everything that \texttt{with}s it recompiles, because
the contract changed. This is automatic and precise: no hand-maintained Makefile
dependency lists, and none of the C ritual where touching one header rebuilds half
the tree.

\cnote{There is no preprocessor and no textual inclusion, so there are no
include-order bugs, no header guards, no macro-hygiene traps, and no
one-definition-rule violations. A unit means the same thing everywhere it is used.}

\textbf{The binder enforces whole-program consistency.} After compilation,
\texttt{gnatbind} checks that every unit in the closure was compiled against
consistent interfaces, then computes a legal \emph{elaboration order} -- the order in
which package-level initialisation runs -- respecting dependencies and any
\texttt{pragma Elaborate\_All}. You cannot accidentally use a package before it has
been elaborated; the binder will not produce a program that does.

\textbf{Projects are the outer layer.} A GNAT project file (\texttt{.gpr}) names the
source directories, the object directory, the switches, and the dependencies on
\emph{other} projects; \texttt{gprbuild} resolves the project graph and the unit
graph together (Chapter~\ref{ch:runtime-build}). In this repository the HAL project
simply globs its \texttt{src/} directory, and GNAT's default file-naming convention
maps a file to a unit by lowercasing the name and turning dots into hyphens:

\begin{lstlisting}[style=sh]
ESP32S3.Stack_Usage   <->   esp32s3-stack_usage.ads / .adb
ESP32S3.W5500.Sockets <->   esp32s3-w5500-sockets.ads / .adb
\end{lstlisting}

\noindent
so dropping the two \texttt{ESP32S3.Stack\_Usage} files into \texttt{src/} made a new
package appear with \emph{no} edit to any build file. An example then needs only

\begin{lstlisting}[style=sh]
with "../../libs/esp32s3_hal/esp32s3_hal.gpr";
\end{lstlisting}

\noindent
and the whole HAL is available to it, every unit named, every dependency tracked.

\section{Packages on bare metal}\index{pragma!Preelaborate}\index{pragma!Pure}

The package is also where the guarantees a restricted runtime cares about attach.
\texttt{pragma Pure} and \texttt{pragma Preelaborate} are package-level promises --
``this unit holds no state'' / ``this unit runs no code at elaboration'' -- and the
compiler \emph{checks} them. On a Zero-Footprint or light-tasking runtime that is not
a nicety: it is what lets a package be used where there is no elaboration machinery
at all. Because a client depends only on a spec, the dependency closure of a
bare-metal application stays small and auditable -- you can enumerate exactly what
gets linked, which feeds straight into the static-footprint accounting of
Chapter~\ref{ch:static-mem} and the \texttt{x mem} tool. And because the spec is the
complete, authoritative interface, the HAL reference (\texttt{./x docs}) is generated
\emph{from the package specs} -- the contract doubles as the documentation.

\section{The unifying idea}

Come back to the three units that opened the chapter -- privacy, separate
compilation, dependency. In C they are three mechanisms, separately maintained, free
to disagree. In Ada they are one: the package, with a compiler-enforced wall between
spec and body. That is why the same construct that hides a hardware register also
bounds a recompilation and also names a dependency edge. This codebase is close to a
pure demonstration of it -- roughly one package per peripheral, a clean
\texttt{ESP32S3.*} tree, a private type guarding every resource, and a build that
rebuilds precisely what changed and nothing else. Once you see the package doing all
three jobs at once, the rest of the architecture in this book reads as the natural
consequence.
